% 2010 年湖南卷理科第 12 题：求 1^2+2^2+...+100^2 的程序框图。
% 来源 PDF：普通高考/2010/2010湖南理.pdf，SHA-256 db6274a80ebb0022ff1b91b2cd3714fe8cca335e48be69805dcd87cdefd8c8d1。
% 原页：第 1 页（A3 三栏）第 2 栏中部，书页码 601。
% 重绘依据：照原页，开始→i=1,s=0→判断 i<=n?；判断框「是」向右、向上依次经
%   s=s+i^2、i=i+1 后向左汇入主干（即在判断前先累加再自增）；「否」向下经输出 s 到结束。
\begin{tikzpicture}[scale=1.0, semithick,
  every node/.style={inner sep=2pt},
  box/.style={draw, minimum width=2.6cm, minimum height=0.6cm, align=center},
  round/.style={draw, rounded corners=6pt, minimum width=1.5cm,
    minimum height=0.6cm, align=center},
  para/.style={draw, trapezium, trapezium left angle=70, trapezium right angle=110,
    minimum width=2.2cm, minimum height=0.65cm, align=center},
  dec/.style={draw, diamond, aspect=2.4, minimum width=2.8cm,
    minimum height=1.1cm, align=center},
  arr/.style={->, >=stealth}]

  \node[round] (start) at (0,9.1) {开始};
  \node[box]   (init)  at (0,8.1) {$i=1,\ s=0$};
  \node[dec]   (test)  at (0,4.9) {$i\leqslant n$?};
  \node[para]  (out)   at (0,3.0) {输出 $s$};
  \node[round] (end)   at (0,1.9) {结束};

  \node[box] (sum) at (3.5,5.9) {$s=s+i^2$};
  \node[box] (inc) at (3.5,7.2) {$i=i+1$};

  \draw[arr] (start) -- (init);
  \draw[arr] (init) -- (test);
  \draw[arr] (test.south) -- (out.north);
  \node[right] at (0.14,3.85) {否};
  \draw[arr] (out) -- (end);

  \draw[arr] (test.east) -| (sum.south);
  \node[right] at (1.75,5.08) {是};
  \draw[arr] (sum.north) -- (inc.south);
  \draw[arr] (inc.west) -- (0,7.2);
\end{tikzpicture}
